% Generated by scripts/materialize_unified_paper.py; do not edit.
\section{First moments, red--blue crossing and asymptotics}
\label{low:sec:rate-lower}

\subsection{The two explicit first-moment thresholds}

Fix \(C\) and \(w>\omega_0\), and let \(\ell\to\infty\).  From \eqref{low:main:eq:2-12},
\eqref{low:main:eq:2-18}, \eqref{low:main:eq:6-1}, and \eqref{low:main:eq:6-12},
\[
\begin{aligned}
\limsup_{\ell\to\infty}\frac1{\ell^2}\log P_{R,\ell}
\le{}&\frac12\log p-\vartheta_C\frac{a^3}{6p^3D}\\
&-G(C,w)-\widehat H(C,w).
\end{aligned}
\tag{L7.1}\label{low:main:eq:7-1}
\]
The discretization term vanishes because \(J_{\ell,d}=o_C(\ell^2)\).
Let \(n=\lfloor e^{\sigma\ell}\rfloor\).  Since
\(\binom n\ell\le n^\ell\), the expected number of red
\(K_\ell\)'s tends to zero whenever
\[
\begin{aligned}
\sigma
&<-\frac12\log p+\vartheta_C\frac{a^3}{6p^3D}
+G(C,w)+\widehat H(C,w)\\
&=-\frac12\log p_C+\frac{B_R(C)}2
+G(C,w)+\widehat H(C,w).
\end{aligned}
\tag{L7.2}\label{low:main:eq:7-2}
\]
This proves the red threshold directly from the probability ledger.

For \(q_\ell=\lfloor C\ell\rfloor\), \eqref{low:main:eq:2-42} and
\(q_\ell/\ell=C+O(\ell^{-1})\) give
\[
 \limsup_{\ell\to\infty}\frac1{\ell^2}\log P_{B,q_\ell}
 \le\frac{C^2}{2}\log(1-p)
 +\frac{2a^3C^3}{3(1-p)^3D}.
\tag{L7.3}\label{low:main:eq:7-3}
\]
The bound \(\binom n{q_\ell}\le n^{q_\ell}\) therefore makes the expected
number of blue \(K_{q_\ell}\)'s tend to zero whenever
\[
 \sigma<-\frac C2\log(1-p)
 -\frac{2a^3C^2}{3(1-p)^3D}
 =-\frac12\log p_C+\frac{CB_B(C)}2.
\tag{L7.4}\label{low:main:eq:7-4}
\]
No blue improvement is used.  A union bound over red and blue vertex sets
therefore produces a coloring with neither forbidden clique for every
\(\sigma\) strictly below the minimum of the right sides of \eqref{low:main:eq:7-2} and
\eqref{low:main:eq:7-4}.

\subsection{The crossing is separated by a constant}

The pinned Gaussian asymptotics, all in the order \(C\to\infty\), are
\[
\begin{gathered}
 \frac p{p_C}=1+O((\log C)^{-1}),\qquad
 \frac{a^3}{p^3D}=\frac12+o(1),\qquad
 \vartheta_C=1-o(1),\\
 B_R(C)=\frac16+o(1),\qquad CB_B(C)=1+o(1).
\end{gathered}
\tag{L7.5}\label{low:main:eq:7-5}
\]
The exact integral \eqref{low:main:eq:2-17} also gives
\[
 G_*(C)=\frac{1+o(1)}{32\log C}=o(1).
\tag{L7.6}\label{low:main:eq:7-6}
\]
Indeed \(m(c-\omega_0)/m_0\to1\), \(P_*(x)\to1\) uniformly, and the
bracket in \eqref{low:main:eq:2-17} tends uniformly to one, while
\(m_0^4/D^2=(1+o(1))/(8\log C)\).

The controlled-residual constants satisfy
\[
\begin{gathered}
 \rho=o(1),\quad m'_0\to1,\quad m_0m''(c)=o(1),\\
 M_w/m_0\to1,\quad j_\pm(w)\to1,\quad\mu_w\to1,
\end{gathered}
\tag{L7.7}\label{low:main:eq:7-7}
\]
where the right limit \(w\downarrow\omega_0\) is taken for each fixed \(C\)
before \(C\to\infty\), and
\[
 \widehat L_w=o(m_0^2),\qquad
 \widehat B_w=m_0^2[1+o(1)].
\tag{L7.8}\label{low:main:eq:7-8}
\]
Thus every positivity and feasibility condition used above holds for all
sufficiently large fixed \(C\), uniformly for \(w\) in a sufficiently small
right-neighborhood of \(\omega_0\).

Equations \eqref{low:main:eq:7-5}--\eqref{low:main:eq:7-8} show
\[
\frac{CB_B(C)}2-
\left[\frac{B_R(C)}2+G_*(C)+\widehat H_*(C)\right]
=\frac5{12}+o(1)>0.
\tag{L7.9}\label{low:main:eq:7-9}
\]
This positive limiting gap, rather than the mere fact that the added terms
are \(o(1)\), proves that red remains the bottleneck.  Taking the supremum
of \eqref{low:main:eq:7-2} over \(w>\omega_0\), using continuity as
\(w\downarrow\omega_0\), and then using \eqref{low:main:eq:7-9} proves
\cref{low:thm:main-lower}.

\subsection{The new leading coefficient}

Finally,
\[
 m_0^2=(2+o(1))\log C,\qquad
 D^2=(32+o(1))(\log C)^3,
\tag{L7.10}\label{low:main:eq:7-10}
\]
and hence
\[
 \widehat B_wA_0=(4+o(1))(\log C)^2,\qquad
 E_w=(2+o(1))(\log C)^2.
\tag{L7.11}\label{low:main:eq:7-11}
\]
Substitution into \eqref{low:main:eq:2-23} gives
\[
 \widehat H_*(C)=
 \frac{(6+o(1))(\log C)^2}
 {12(32+o(1))(\log C)^3}
 =\frac{1+o(1)}{64\log C}.
\tag{L7.12}\label{low:main:eq:7-12}
\]
The exact-merge predecessor is
\(H_*(C)=(1+o(1))/(96\log C)\), so the new term is eventually strictly
larger.  Since \(B_R(C)/2\to1/12\), the public coarsening
\(B_R(C)/2\ge1/20\) also follows after increasing the existential threshold.

\subsection{A precisely delimited scalar method cap}

\begin{definition}[Admissible scalar residual family]\label{low:def:admissible}
Fix constants \(M_0,M_{\rm ext}>0\).  An
\((M_0,M_{\rm ext})\)-admissible family \(\mathcal E\) assigns, for every
sufficiently large \(C\), every \(w\) in a right-neighborhood of
\(\omega_0\), and every \(0\le r\le\ell\), deterministic data
\[
 (e^{C,w,\ell,r},T^{C,w,\ell,r},
   \beta_{C,w,\ell,r},\mu_{C,w,\ell,r}).
\]
The same constants \(M_0,M_{\rm ext}\) apply to every member of the family
and to every \(C,w,\ell,r\).  At \(r=0,1\), the arrays are empty and we set
\(\beta=0\), \(\mu=1\), and \(\mathcal G_{\mathcal E}=0\).
\begin{enumerate}[label=(A\arabic*)]
\item The increments are fixed before exposure and
\[
 0\le e_{ij}\le M_0m_0^2r\qquad(1\le i<j\le r).
\tag{L7.13}\label{low:main:eq:7-13}
\]
\item The array is the unique solution of the generalized recursion
\[
 T_{ij}=t_0+\gamma_d
 \sum_{\substack{q>i\\q\ne j}}T_{jq}+e_{ij},
 \qquad \Delta_{ij}=T_{ij}-t_0\ge0.
\tag{L7.14}\label{low:main:eq:7-14}
\]
Thus neither the recursion nor the admissible class refers to the particular
choice \(e_{ij}=\theta(i-1)\).
\item Put
\[
 S_i=\sum_{\{j,q\}\subset V_i}\Delta_{jq},\qquad
 R_i=\max_{j\in V_i}\sum_{q\ne j}\Delta_{jq},\qquad k=|V_i|.
\]
The three-factor split is feasible at every step:
\[
 \frac{4t_0k}{d}+\frac{4R_i}{d}\le\frac12.
\tag{L7.15}\label{low:main:eq:7-15}
\]
For \(r\ge2\), writing \(F_i^T\) for \eqref{low:main:eq:3-1} with coefficient array \(T\), the
two scalar coefficients obey the exact constraints
\[
 0<\mu_{C,w,\ell,r}\le
 \inf_{1\le i<r}\ \inf_{b\in[c-w,c+w]^{V_i}}
 \lambda_{\min}\!\bigl(-\nabla^2F_i^T(b)\bigr),
 \qquad 0\le\beta_{C,w,\ell,r}\le m_0^2.
\tag{L7.16}\label{low:main:eq:7-16}
\]
With these coefficients, define
\[
 \lambda_i^{\mathcal E}=\mathsf d_{r-i,w,d}
 +\beta_{C,w,\ell,r}\frac{S_i}{d}
 -\sum_{j>i}\frac{e_{ij}^2}{2\mu_{C,w,\ell,r}d}
 -\frac{d^{-1/5}}{\sqrt d}\sum_{j>i}T_{ij}.
\tag{L7.17}\label{low:main:eq:7-17}
\]
The family must certify, on every history in
\(\mathfrak H^R_{r,i-1}\), the one-step inequality
\[
 \mathcal K_i^T\le
 \frac{\mathcal H_C(r-i+1)}{\mathcal H_C(r-i)}
 e^{-\lambda_i^{\mathcal E}}.
\tag{L7.18}\label{low:main:eq:7-18}
\]
Thus this scalar class credits exactly the displayed \(\beta S_i/d\) term,
pays the displayed scalar square penalty, and discards every favorable term
not listed in \eqref{low:main:eq:7-17}.
\item Relative to the weighted propagation lemma (S3), the family's credited
deterministic contribution is the normalized functional
\[
 \mathcal G_{\mathcal E}(C,w,\ell,r)
 =\frac1d\left[
 \beta_{C,w,\ell,r}\sum_{i=1}^{r-1}S_i
 -\frac1{2\mu_{C,w,\ell,r}}
  \sum_{1\le i<j\le r}e_{ij}^2\right].
\tag{L7.19}\label{low:main:eq:7-19}
\]
This definition excludes adaptive weights, matrix-valued completion, and
unlisted source credits.
\item After deleting \(u\) vertices, relabeling in retained order, and
rebuilding the family member, with \(x_+=\max\{x,0\}\),
\[
 0\le
 \left[
 \mathcal G_{\mathcal E}(C,w,\ell,\ell)
 -\mathcal G_{\mathcal E}(C,w,\ell,\ell-u)
 \right]_+
 \le\frac{M_{\rm ext}}{\log C}\,\ell u
\tag{L7.20}\label{low:main:eq:7-20}
\]
for every \(0\le u\le\ell\).  Thus extraction is a quantified condition,
not a reference to the optimizer-specific estimate.
\end{enumerate}
Define
\[
\begin{aligned}
 \operatorname{Lead}_C(\mathcal E)
 &=\limsup_{w\downarrow\omega_0}\,
   \limsup_{\ell\to\infty}
   \frac{\mathcal G_{\mathcal E}(C,w,\ell,\ell)}{\ell^2},\\
 \mathfrak M_{M_0,M_{\rm ext}}(C)
 &=\sup_{\mathcal E}\operatorname{Lead}_C(\mathcal E),
\end{aligned}
\tag{L7.21}\label{low:main:eq:7-21}
\]
where the supremum is over \((M_0,M_{\rm ext})\)-admissible families.  The
right-neighborhood in the first limsup may depend on \(C\), but not on
\(\ell\) or \(r\).
\end{definition}

\begin{proposition}[Cap within \cref{low:def:admissible}]\label{low:prop:method-cap}
For every fixed \(M_0,M_{\rm ext}\), the leading functional in \eqref{low:main:eq:7-21}
satisfies
\[
 \limsup_{C\to\infty}64\log C\,
 \mathfrak M_{M_0,M_{\rm ext}}(C)\le1.
\tag{L7.22}\label{low:main:eq:7-22}
\]
There are fixed absolute constants \(M_0^*,M_{\rm ext}^*\) for which the
family \(e_{ij}=\theta(i-1)\), \(\beta=\widehat B_w(C)\), and
\(\mu=\mu_w\) is admissible and satisfies
\[
 \lim_{C\to\infty}64\log C\,
 \operatorname{Lead}_C(\mathcal E^*)=1.
\tag{L7.23}\label{low:main:eq:7-23}
\]
Thus \(1/64\) is both the cap and an attained leading coefficient within the
precisely delimited scalar class, for fixed admissibility constants large
enough to contain \(\mathcal E^*\).
\end{proposition}

\begin{proof}
Fix an admissible family and abbreviate its two scalar coefficients by
\(\beta\) and \(\mu\).  The exact constraint \eqref{low:main:eq:7-16} gives
\[
 \beta\le m_0^2.
\tag{L7.24}\label{low:main:eq:7-24}
\]
At \(b=c\mathbf1\), the row condition \eqref{low:main:eq:7-15} and the diagonal formula \eqref{low:app-A:eq:a-1}
give, for every \(j\in V_i\),
\[
 \mu\le\lambda_{\min}(-\nabla^2F_i^T(c\mathbf1))
 \le(-\nabla^2F_i^T(c\mathbf1))_{jj}
 \le m'_0+\frac18m_0m''(c)=1+o(1).
\tag{L7.25}\label{low:main:eq:7-25}
\]
Here \(t_0k+R_i\le d/8\), \(m'_0\to1\), and
\(m_0m''(c)=o(1)\).  Hence every valid common scalar completion constant
satisfies, uniformly over the class,
\[
 \mu\le1+o(1).
\tag{L7.26}\label{low:main:eq:7-26}
\]

Let \(\mathcal L\) be the feedback operator in \eqref{low:main:eq:7-14}.  Its row norm is at
most \(\gamma_dr\le m_0m'_0/D=\rho=o(1)\).  Writing
\[
 g_{ij}=A_0(r-i-1),
\]
the recursion for \(\Delta\) is
\(\Delta=g+\mathcal L\Delta+e\).  Therefore
\[
 \Delta=(I-\mathcal L)^{-1}g+e+O(\rho m_0^2r)
\tag{L7.27}\label{low:main:eq:7-27}
\]
uniformly on every edge, with an implicit constant depending only on \(M_0\),
by \eqref{low:main:eq:7-13}.  The exact exposure multiplicity is
\[
 \sum_{i=1}^{r-1}S_i
 =\sum_{1\le j<q\le r}(j-1)\Delta_{jq}.
\tag{L7.28}\label{low:main:eq:7-28}
\]
The first term of \eqref{low:main:eq:7-27} contributes
\[
 2A_0\binom r4[1+O(\rho)]
\tag{L7.29}\label{low:main:eq:7-29}
\]
to \eqref{low:main:eq:7-28}.  The last term changes the numerator in \eqref{low:main:eq:7-19} by at most
\(O_{M_0}(\rho m_0^4r^4)=o(m_0^4r^4)\).

For the direct residual term, scalar square completion is pointwise:
\[
 \beta(j-1)e_{jq}-\frac{e_{jq}^2}{2\mu}
 \le\frac{\mu\beta^2}{2}(j-1)^2.
\tag{L7.30}\label{low:main:eq:7-30}
\]
Summing \eqref{low:main:eq:7-30} gives \(\mu\beta^2K_r/2\).  Set \(r=\ell\), divide \eqref{low:main:eq:7-19}
by \(\ell^2\), and use
\(\binom\ell4/\ell^4\to1/24\), \(K_\ell/\ell^4\to1/12\), and
\(d/\ell^2\to D^2\).  Equations \eqref{low:main:eq:7-24}--\eqref{low:main:eq:7-30} yield
\[
 \operatorname{Lead}_C(\mathcal E)
 \le\frac{\beta A_0}{12D^2}
   +\frac{\mu\beta^2}{24D^2}
   +o\!\left(\frac{m_0^4}{D^2}\right)
 \le\frac{m_0^4}{8D^2}[1+o(1)]
 =\frac{1+o(1)}{64\log C}.
\tag{L7.31}\label{low:main:eq:7-31}
\]
The error is uniform for fixed \(M_0,M_{\rm ext}\), so taking the supremum
and then the limsup proves \eqref{low:main:eq:7-22}.

For \(e_{ij}=\theta(i-1)\), equations \eqref{low:main:eq:5-5}, \eqref{low:main:eq:7-7}, and \eqref{low:main:eq:7-8} make
\eqref{low:main:eq:7-30} asymptotically tight.  Equations \eqref{low:main:eq:5-11}, \eqref{low:main:eq:6-2}, and \eqref{low:main:eq:4-4} verify
\eqref{low:main:eq:7-15}--\eqref{low:main:eq:7-18}, and \eqref{low:main:eq:7-13} holds with \(M_0^*=1\).  To verify \eqref{low:main:eq:7-20}, write
\(\Delta^{(r)}\) for the array rebuilt at target \(r\), and put
\(r_-=\ell-u\).  Subtracting the two positive recursions \eqref{low:main:eq:7-14}, and using
\(\|\mathcal L\|_{\infty}\le\rho<1/2\), gives absolute constants
\(K_1,K_2\) such that
\[
 \max_{1\le j<s\le r_-}
 \left|\Delta_{js}^{(\ell)}-\Delta_{js}^{(r_-)}\right|
 \le K_1m_0^2u,
 \qquad
 \max_{j<s\le\ell}\Delta_{js}^{(\ell)}\le K_2m_0^2\ell.
\tag{L7.32}\label{low:main:eq:7-32}
\]
Indeed, the second bound is \eqref{low:main:eq:5-1}.  If \(U\) denotes the first maximum,
subtraction of \eqref{low:main:eq:7-14} gives
\(U\le A_0u+\rho U+\gamma_duK_2m_0^2\ell\).  Now
\(A_0=O(m_0^2)\), \(\gamma_d\ell\le\rho\), and \(\rho<1/2\) give the first
bound.  The multiplicity identity \eqref{low:main:eq:7-28} and \eqref{low:main:eq:5-4} now give
\[
 \left|\sum_iS_i^{(\ell)}-\sum_iS_i^{(r_-)}\right|
 \le K_3m_0^2\ell^3u,
 \qquad |K_\ell-K_{r_-}|\le K_4\ell^3u
\tag{L7.33}\label{low:main:eq:7-33}
\]
for absolute \(K_3,K_4\).  Since \(\beta\le m_0^2\), \(\mu_w\to1\), and
\(m_0^4/D^2=O(1/\log C)\), \eqref{low:main:eq:7-33} proves \eqref{low:main:eq:7-20} for some absolute
\(M_{\rm ext}^*\).  Finally, the functional is exactly \eqref{low:main:eq:2-23} up to
\(o(1/\log C)\), and \eqref{low:main:eq:7-12} gives the asserted limit \eqref{low:main:eq:7-23}.
\end{proof}

The proposition does not quantify over outcome-adaptive weights,
matrix-valued square completion, higher-order potentials, a sharper source
quadratic MGF, or a different extraction invariant.
